\chapter{Orchitis}
\index{Orchitis}

\medskip
\noindent\textit{\textbf{Under review.} This chapter is an AI-assisted educational draft; it is not a substitute for professional medical advice.}

\section*{Definition and Overview}
inflammation of one or both testes, often associated with viral infection or bacterial epididymo-orchitis.

\section*{Clinical Features}
testicular pain and swelling, fever, nausea, urinary symptoms, or discharge. Sudden severe unilateral pain requires urgent exclusion of torsion.

\section*{When to Seek Medical Care}
Seek medical review for new, persistent, worsening, or function-limiting symptoms. Contact your local emergency services for severe breathing difficulty, collapse, sudden neurological change, severe chest or abdominal pain, or any rapidly worsening illness.

\section*{Diagnosis and Investigations}
examination, urine and STI tests where indicated, blood tests, and scrotal ultrasound when torsion, abscess, or another acute cause is possible.

\section*{Management}
analgesia, scrotal support, rest, and antibiotics when bacterial infection is suspected. Testicular torsion is an emergency requiring surgery.

\section*{Complications}
Complications depend on the underlying cause and may include progression, effects on other organs, treatment adverse effects, or reduced quality of life. Early assessment can reduce preventable harm.

\section*{Prognosis and Self-Care}
Outcome depends on the cause, severity, complications, and response to treatment. Follow the care plan, keep appointments, avoid known triggers, and seek help if symptoms worsen.
\section*{Safety-netting}
Seek urgent care for sepsis features, breathing difficulty, confusion, severe dehydration, meningitis symptoms, or rapid deterioration.
